\chapter{การทดลอง ประเมินผล และอภิปรายผล}

\section{ผลการประมาณค่า $p(e)$ จาก CVSS (วัตถุประสงค์ข้อ 2)}

edge ที่เป็น misconfiguration ไม่มี CVE จึงไม่มีคะแนน EPSS ให้ใช้โดยตรง
โครงงานนี้จึงกำหนดเวกเตอร์ CVSS ให้แต่ละคลาสของ misconfiguration ตาม rubric (ภาคผนวก ก)
แล้วใช้ฟังก์ชันที่เรียนรู้ความสัมพันธ์ระหว่างเวกเตอร์ CVSS กับคะแนน EPSS ของ CVE จริง
แปลงเวกเตอร์นั้นเป็นค่าความน่าจะเป็น $p(e)$
หัวข้อนี้รายงานว่าฟังก์ชันดังกล่าวแม่นเพียงใด
เพราะความคลาดเคลื่อนของ $p(e)$ ส่งผลต่ออันดับเส้นทางโดยตรง

\subsection{ข้อมูลและวิธีการ}

ข้อมูลที่ใช้เป็น CVE จำนวน 46{,}360 รายการที่มีทั้งเวกเตอร์ CVSS (จาก NVD จำนวน 33 หน้า)
และคะแนน EPSS จาก snapshot วันที่ 2026-08-26 โดยเป็น CVE ที่เผยแพร่ระหว่าง 2017-03-20 ถึง 2026-08-06
การสุ่มทุกขั้นตอนใช้ seed 20260826
ตัวแปรตามคือ EPSS ในสเกล logit
และเปรียบเทียบ 4 โมเดล ได้แก่ ค่าเฉลี่ยคงที่ (constant baseline),
การถดถอยเชิงเส้นจากคะแนนฐาน CVSS เพียงตัวเดียว,
การถดถอยเชิงเส้นจาก 8 metric ของ CVSS (OLS)
และ gradient boosting จาก 8 metric เดียวกัน (GBM)
ตัวชี้วัดความแม่นในตารางที่~\ref{tab:fit} คำนวณบนข้อมูลทดสอบ 25\%

\subsection{ผลการ fit}

\begin{table}[H]
\centering
\caption{ผลการ fit จากเวกเตอร์ CVSS ไปเป็น EPSS (ข้อมูลทดสอบ 25\%)}\label{tab:fit}
\footnotesize\setlength{\tabcolsep}{4pt}
\begin{tabular}{p{5.6cm}rrrr}
\toprule
โมเดล & $R^2$ (logit) & $R^2$ (prob) & $R^2$ (top 10\%) & $\rho$ \\
\midrule
constant baseline      & 0.000 & $-0.032$ & $-0.452$ & --- \\
OLS จากคะแนนฐานตัวเดียว & 0.118 & $-0.018$ & $-0.432$ & 0.354 \\
\textbf{OLS จาก 8 metric (ใช้จริง)} & \textbf{0.219} & $-0.010$ & $-0.420$ & 0.548 \\
GBM จาก 8 metric (ไม่ได้ใช้) & 0.268 & $-0.005$ & $-0.415$ & 0.601 \\
\bottomrule
\end{tabular}
\end{table}

ผลลัพธ์สำคัญมี 3 ข้อ
ข้อแรก โมเดลที่ใช้ตั้งราคา edge คือ OLS จาก 8 metric
เพราะสัมประสิทธิ์แต่ละตัวบวกกันในสเกล logit จึงตรวจสอบย้อนกลับได้ทีละ metric
ซึ่งตรงกับวัตถุประสงค์ว่า $p(e)$ ต้องอธิบายที่มาได้
GBM ให้ค่าสูงกว่าเล็กน้อยแต่เป็นกล่องดำและไม่ได้นำไปใช้ ตัวเลขทั้งหมดในรายงานจึงเป็นของ OLS
ข้อสอง $R^2$ ในสเกลความน่าจะเป็นติดลบในทุกโมเดล
หมายความว่าในหน่วยที่ $p(e)$ ถูกใช้จริง ฟังก์ชันที่ fit ได้ไม่ดีกว่าการทายค่าเฉลี่ย
และอธิบายความแปรปรวนได้เพียงประมาณหนึ่งในห้า (0.219) เฉพาะในสเกล logit
ข้อสาม ความคลาดเคลื่อนไม่ได้เป็นสัญญาณรบกวนเท่ากันทุกช่วง
แต่เกิดจากการที่ฟังก์ชัน ``ไปไม่ถึงหาง'' ของการแจกแจง
ค่าที่ทำนายสูงสุดคือ 0.0186 ขณะที่ EPSS จริงสูงสุดถึง 0.99999
ฟังก์ชันจึงถดถอยเข้าหาค่าฐาน (base rate) ประมาณ 0.0055

ค่า $R^2 = 0.219$ ในตารางที่~\ref{tab:fit} ใช้การแบ่งข้อมูลทดสอบแบบสุ่มราย CVE
ซึ่งมองโลกในแง่ดีกว่าการใช้งานจริงเล็กน้อย เพราะ CVE ที่คล้ายกันอาจตกอยู่คนละฝั่งของการแบ่ง
เมื่อกันทั้งหน้า NVD ออกจากกัน $R^2$ ลดเหลือ 0.156 และเมื่อกันทั้งปีออกเหลือ 0.150
กรณีที่เข้มที่สุดคือฝึกด้วย CVE ถึงปี 2022 แล้วทดสอบกับปี 2023 เป็นต้นไป $R^2$ ติดลบ ($-0.863$)
แต่ Spearman $\rho$ ยังคงอยู่ที่ 0.542 แปลว่า \emph{ระดับ} ของ $p(e)$ เลื่อนไปตามเวลา
ขณะที่ \emph{อันดับ} ยังรักษาไว้ได้ (ตรวจด้วย \texttt{scripts/audit\_data.py})
จึงไม่ควรอ้างว่าฟังก์ชันนี้ทำนายค่า EPSS ของ CVE ใหม่ได้ ใช้ได้เพียงเพื่อจัดอันดับ

\subsection{ความคลาดเคลื่อนในหน่วย log}

เนื่องจากต้นทุนของเส้นทางคือผลรวมของ $-\log p$ บน edge
ค่าที่กำหนดว่าอันดับเส้นทางเปลี่ยนหรือไม่คือการกระจายของ
$r = \log p_{\text{จริง}} - \log p_{\text{ทำนาย}}$
ไม่ใช่ค่าสหสัมพันธ์ $\rho$ เพราะการแปลงแบบ monotone ต่อ edge สามารถสลับอันดับของเส้นทางได้
(เส้นทางที่มี 2 edge ที่ $p=0.9$ ได้ $0.81$ ชนะเส้นทางที่มี 1 edge ที่ $p=0.8$
แต่ถ้า map $0.9 \mapsto 0.5$ และ $0.8 \mapsto 0.45$ อันดับจะกลับกัน)
อันดับของเส้นทางคงเดิมพอดีเมื่อฟังก์ชันเป็น power law $\hat{p} \propto p^{k}$
จึงนำค่ามัธยฐานของ $r$ ออก (ซึ่งคือค่า $k$) แล้ววัดความกว้างที่เหลือ
ตารางที่~\ref{tab:resid} รายงานค่าสัมบูรณ์ $|r|$ จากการทำนายแบบ 5-fold out-of-fold บน CVE ทั้ง 46{,}360 รายการ
โดยแยกกลุ่มที่ EPSS สูงที่สุด 10\% ออกมาต่างหาก เพราะ edge ที่ตัดสินเส้นทางคือ edge ที่มีโอกาสสูง

\begin{table}[H]
\centering
\caption{ความคลาดเคลื่อน $|r|$ ในหน่วย nats แยกตามช่วงของ EPSS จริง}\label{tab:resid}
\small
\begin{tabular}{lrrr}
\toprule
กลุ่ม & $n$ & $|r|$ มัธยฐาน & $|r|$ เปอร์เซ็นไทล์ที่ 90 \\
\midrule
ทั้งหมด (pooled) & 46{,}360 & 0.587 & 1.483 \\
ล่าง 90\% & 41{,}724 & 0.523 & 1.172 \\
\textbf{EPSS สูงสุด 10\%} & 4{,}636 & \textbf{1.938} & \textbf{3.980} \\
\bottomrule
\end{tabular}
\end{table}

กลุ่มที่ EPSS สูงสุด 10\% คลาดเคลื่อนมากกว่ากลุ่มล่าง 90\% ประมาณ 3.7 เท่า (1.938 เทียบกับ 0.523)
และการรายงานค่า pooled เพียงค่าเดียว (0.587) จะทำให้ฟังก์ชันดูแม่นเกินจริง
ค่า 1.938 nats เทียบเท่าตัวคูณ $e^{1.938} \approx 6.95$ เท่าบน $p(e)$
ขณะที่แถบการรบกวน $\pm 10$--$50\%$ ในข้อเสนอโครงงานมีขนาดมากสุดเพียง $\log 1.5 \approx 0.405$ nats (1.5 เท่า)
แถบที่กำหนดไว้จึงแคบกว่าที่วัดได้ประมาณ 4.8 เท่าที่มัธยฐาน และ 9.8 เท่าที่เปอร์เซ็นไทล์ที่ 90
(ในหน่วย log) ข้อสรุปที่ถูกต้องจึงเป็น ``แคบกว่า 4--10 เท่า แล้วแต่เปอร์เซ็นไทล์''
ไม่ใช่ ``หนึ่ง order of magnitude'' ซึ่งเป็นจริงเฉพาะที่เปอร์เซ็นไทล์ที่ 90

\subsection{ผลต่อความสามารถในการแยกอันดับเส้นทาง}

เมื่อ $|r|\le\delta$ ทุก edge เส้นทางสองเส้นที่มีความยาว $|P|$ และ $|Q|$ จะรักษาอันดับจริงไว้ได้
ก็ต่อเมื่อผลต่างต้นทุนจริงมากกว่า $\delta\,(|P|+|Q|)$
ตารางที่~\ref{tab:order} แสดงขั้นต่ำของผลต่างนั้น เมื่อใช้ $\delta$ เป็นมัธยฐานและเปอร์เซ็นไทล์ที่ 90

\begin{table}[H]
\centering
\caption{ผลต่างต้นทุนขั้นต่ำ (nats) ที่ทำให้สองเส้นทางยาวเท่ากันยังรักษาอันดับไว้ได้}\label{tab:order}
\small
\begin{tabular}{rrrrr}
\toprule
& \multicolumn{2}{c}{ค่า pooled} & \multicolumn{2}{c}{EPSS สูงสุด 10\%} \\
\cmidrule(lr){2-3}\cmidrule(lr){4-5}
ความยาวเส้นทาง & มัธยฐาน & P90 & มัธยฐาน & P90 \\
\midrule
2 edge & 2.35 & 5.93 & 7.75 & 15.92 \\
4 edge & 4.70 & 11.86 & 15.51 & 31.84 \\
8 edge & 9.39 & 23.72 & 31.01 & 63.68 \\
\bottomrule
\end{tabular}
\end{table}

สำหรับเส้นทางยาว 4 edge ผลต่างต้นทุน 4.70 nats เทียบเท่าความน่าจะเป็นจริงที่ต่างกัน
$e^{4.70} \approx 109$ เท่า (edge คุณภาพมัธยฐาน)
และ 15.51 nats ในกลุ่มหาง เทียบเท่าประมาณ 5.4 ล้านเท่า
ดังนั้นน้ำหนักจาก fit แยกอันดับเส้นทางได้เฉพาะคู่ที่ต่างกันมากเท่านั้น
และแทบแยกไม่ได้เลยเมื่อ edge ที่ตัดสินอยู่ในกลุ่มที่ EPSS สูง
เงื่อนไขนี้เป็นเงื่อนไขจำเป็น ไม่ใช่เงื่อนไขเพียงพอ
ผลที่เกิดจริงต่ออันดับจุดตัดถูกวัดโดยตรงในหัวข้อ~\ref{sec:exp6}

\section{การเปรียบเทียบ Dominator Tree กับ Weighted Min-Cut (วัตถุประสงค์ข้อ 3)}

วัตถุประสงค์ข้อ 3 ถามว่าจุดตัดสองนิยามให้คำตอบตรงกันหรือไม่ นิยามแรกคือ post-dominator
ซึ่งเป็นโครงสร้างล้วน ไม่อ่านน้ำหนัก และคืนเซตของ node ที่อยู่บน\emph{ทุก}เส้นทาง
นิยามที่สองคือ weighted node min-cut ซึ่งคืนชุด node ที่ตัดแล้วผู้โจมตีไปไม่ถึงเป้าหมาย
โดยมีต้นทุนรวมต่ำสุด การเทียบสองนิยามนี้ทำบนกราฟชุดเดียวกันทั้ง 864 กราฟถ่วงน้ำหนัก

\subsection{สามนิยามของจุดตัด}

ก่อนเทียบกับ min-cut ต้องแยกให้ชัดว่า ``จุดตัด'' หมายถึงนิยามใด เพราะทั้งสามให้เซตที่ต่างกัน

\begin{table}[H]
\centering
\caption{สามนิยามของจุดตัดบน 864 กราฟถ่วงน้ำหนัก}\label{tab:defs}
\small
\begin{tabular}{p{4.6cm}rr}
\toprule
นิยาม & กราฟที่เซตไม่ว่าง & ขนาดเซตเฉลี่ย \\
\midrule
global (ทุก entry ถึงทุก target) & 8.3\% & 4.8 \\
per-target (ทุก entry ถึง target เดียว) & 100\% & 5.8 \\
per-pair (หนึ่งคู่ entry--target) & 100\% & 9.2 \\
\bottomrule
\end{tabular}
\end{table}

นิยาม global เข้มที่สุด จุดตัดต้องอยู่บนทุกเส้นทางของทุกคู่ จึงพบเพียง 18 จาก 216 โครงสร้าง (8.3\%)
ส่วน per-target และ per-pair ไม่ว่าง 100\% ซึ่ง \textbf{เป็น artifact ของตัวสร้างกราฟ ไม่ใช่ผลต่อเครือข่ายจริง}
ในกราฟที่สร้างขึ้น target แต่ละตัวมี in-edge เพียงเส้นเดียวโดยโครงสร้าง node ต้นทางของ edge นั้นจึง
ครอบทุกเส้นทางที่เข้าถึง target เสมอ ค่า 100\% จึงสะท้อนวิธีสร้างกราฟ ไม่ใช่คุณสมบัติของเครือข่ายจริง
รายงานนี้จึงใช้นิยาม global เป็นพาดหัว และกำกับข้อจำกัดนี้ไว้ทุกที่ที่อ้างถึง per-target/per-pair

\subsection{จุดตัดเชิงโครงสร้างเป็นกรณีสุดขั้วของชั้นความน่าจะเป็น}

post-dominator set พิสูจน์ได้ว่าเท่ากับ $\{v : \Delta(v) = \infty\}$ นั่นคือ node ที่ลบแล้ว
ผู้โจมตีไปไม่ถึงเป้าหมายเลย (ต้นทุนเพิ่มเป็นอนันต์) จุดตัดเชิงโครงสร้างจึงไม่ใช่วิธีที่แยกจาก
ชั้นความน่าจะเป็น แต่เป็นกรณีสุดขั้วของมัน คุณสมบัตินี้ตรวจด้วย brute force บนกราฟสุ่ม 300 กราฟ
(\texttt{tests/test\_bruteforce.py}): จำนวน node ที่ $\Delta(v)=\infty$ เท่ากับขนาดของ global set ทุกกราฟ

\subsection{ความตรงกันของ dominator กับ min-cut}

การเทียบทำบนกราฟ 18 โครงสร้างที่ global set ไม่ว่าง (เซตว่างเทียบไม่ได้ และ \texttt{disagreement}
บนเซตว่างมีค่า 1 โดยปริยาย ไม่ได้บอกอะไรเกี่ยวกับสองนิยาม) min-cut รายงานสาม capacity ตามการตัดสินใจ
ตารางที่~\ref{tab:cut} แสดงสัดส่วนที่ node ใน min-cut เป็นสมาชิกของ dominator set ด้วย

\begin{table}[H]
\centering
\caption{สัดส่วน node ใน min-cut ที่อยู่ใน dominator set (18 โครงสร้างที่ global ไม่ว่าง)}\label{tab:cut}
\small
\begin{tabular}{lrp{3.2cm}l}
\toprule
capacity & ขนาด cut เฉลี่ย & cut อยู่ใน dominator set & บทบาท \\
\midrule
unit & 1.0 & 100\% & หน่วยตรงกับ dominator \\
prob & 1.1 & 72\% EPSS, 100\% อื่น & ค่าตั้งต้น \\
cost & 3.0 & 10\% & degenerate \\
\bottomrule
\end{tabular}
\end{table}

ผลสำคัญมีสามข้อ ข้อแรก capacity แบบ \texttt{unit} ซึ่งนับจำนวน node ที่ต้องตัด ตรงกับ dominator set
100\% เป็นผลที่ควรเป็น เพราะทั้งคู่วัดในหน่วยเดียวกันและ dominator ไม่อ่านน้ำหนัก
ข้อสอง capacity แบบ \texttt{prob} ซึ่งเป็นค่าตั้งต้น ตรงกับ dominator 100\% เมื่อใช้น้ำหนัก CVSS หรือ uniform
แต่ลดเหลือ \textbf{72\% เมื่อใช้น้ำหนัก EPSS} กล่าวคือใน 28\% ของกรณี min-cut ที่ถ่วงด้วย EPSS
เลือก node ที่ไม่ใช่ dominator เพราะ node นั้นเป็นคอขวดของ\emph{ความน่าจะเป็น} แม้ไม่ได้อยู่บนทุกเส้นทาง
นี่คือข้อค้นพบหลักของวัตถุประสงค์ข้อ 3 ที่น้ำหนักเชิงประจักษ์ให้คำตอบต่างจากโครงสร้างล้วน
ข้อสาม capacity แบบ \texttt{cost} ($-\log p$) เสื่อม (degenerate) เพราะ edge เชิงโครงสร้างมี $p=1$ จึง $w=0$
node ที่เข้าถึงได้ฟรีจึงถูกตัดด้วยต้นทุนศูนย์ min-cut ยุบไปที่เซตต้นทุนศูนย์ที่เลือกโดยพลการ
ยืนยันด้วยสัดส่วน cut ที่มีต้นทุนรวมเป็นศูนย์พอดีถึง 77\% ของแถว จึงรายงาน \texttt{cost} ไว้เป็น
คอลัมน์ sensitivity เท่านั้น ไม่ใช่ค่าหลัก

\subsection{ระดับการแบ่งส่วนเครือข่ายกับการเกิดจุดตัด}

ข้อเสนอโครงงานถามว่าเครือข่ายต้องแบ่งส่วนละเอียดแค่ไหนจุดตัดที่การันตีได้จึงจะปรากฏ
คำถามนี้ตอบได้จากสัดส่วนที่ global set ไม่ว่าง แยกตามระดับการแบ่งส่วน รายงานเป็นจำนวนกราฟ
พร้อมช่วงความเชื่อมั่น 95\% แบบ Wilson (ไม่ใช่เปอร์เซ็นต์ลอย ๆ เพราะจำนวนน้อย)

\begin{table}[H]
\centering
\caption{สัดส่วนกราฟที่มี global choke point แยกตามการแบ่งส่วน (72 กราฟต่อระดับ)}\label{tab:seg}
\small
\begin{tabular}{lrl}
\toprule
การแบ่งส่วน & กราฟที่มีจุดตัด & ช่วงความเชื่อมั่น 95\% (Wilson) \\
\midrule
flat (ไม่แบ่ง) & 1 / 72 & 0.2--7.5\% \\
3-tier & 1 / 72 & 0.2--7.5\% \\
micro-segmentation & 16 / 72 & 14--33\% \\
\midrule
รวม & 18 / 216 & 5.3--12.8\% \\
\bottomrule
\end{tabular}
\end{table}

micro-segmentation ให้จุดตัดที่การันตีได้บ่อยกว่าอีกสองแบบอย่างมีนัยสำคัญ (ช่วงความเชื่อมั่นไม่ซ้อนทับ)
ซึ่งสอดคล้องกับสัญชาตญาณ เครือข่ายที่แบ่งส่วนละเอียดบีบให้เส้นทางลอดผ่าน node เชื่อมไม่กี่ตัว
node เหล่านั้นจึงกลายเป็นคอขวดที่ทุกเส้นทางต้องผ่าน ส่วน flat และ 3-tier มีเส้นทางสำรองมากเกินกว่าจะมี
จุดตัดร่วม อย่างไรก็ตามแม้แต่ micro ก็ยังมีจุดตัดเพียง 22\% ของกราฟ แปลว่า\textbf{การพึ่งจุดตัดที่การันตีได้
เป็นกลยุทธ์ที่ใช้ได้เฉพาะส่วนน้อยของกราฟ} ส่วนที่เหลือต้องใช้การจัดอันดับ $\Delta(v)$ แทนการหาเซตที่การันตี

หมายเหตุด้านความแปรปรวน: grid รุ่นก่อน (135 โครงสร้าง, seed คนละชุด) เคยให้ค่า micro ที่ 44\%
ค่าที่ต่างกันมากนี้ย้ำว่าสัดส่วนดังกล่าวมีความแปรปรวนสูงเมื่อจำนวนกราฟที่เข้าเงื่อนไขมีน้อย
จึงต้องรายงานเป็นจำนวนกราฟพร้อมช่วงความเชื่อมั่นเสมอ ไม่ใช่ตัวเลขจุดเดียว
\section{ความเสถียรของอันดับเมื่อน้ำหนักถูกรบกวน (วัตถุประสงค์ข้อ 4)}

วัตถุประสงค์ข้อ 4 ถามว่าอันดับจุดตัดขยับแค่ไหนเมื่อน้ำหนัก edge คลาดเคลื่อน การทดลองนี้ใช้สอง
รูปแบบของสัญญาณรบกวน (noise) คือแบบสมมาตร $\pm 10$--$50\%$ ตามที่ข้อเสนอโครงงานระบุ
และแบบเชิงประจักษ์ (empirical) ซึ่งสุ่มจาก residual จริงของการ fit ในหัวข้อ~\ref{sec:exp6}
คูณด้วยสเกล 0.5 ถึง 8 เท่า (สเกล $\geq 4$ เพิ่มเข้ามาเพื่อหาจุดที่อันดับเริ่มพัง เกินกว่าที่ข้อเสนอกำหนด)
วัดด้วย Jaccard@5 เทียบอันดับก่อนและหลังรบกวน บนกราฟเดิมและโมเดลน้ำหนักเดิม

\begin{table}[H]
\centering
\caption{Jaccard@5 เฉลี่ยเมื่อรบกวนน้ำหนัก (ทุกโมเดลน้ำหนักรวมกัน)}\label{tab:noise}
\small
\begin{tabular}{lrrrrr}
\toprule
& \multicolumn{4}{c}{แบบสมมาตร ($\pm$)} & \\
\cmidrule(lr){2-5}
& 10\% & 20\% & 30\% & 50\% & \\
Jaccard@5 & 0.979 & 0.968 & 0.964 & 0.954 & \\
\midrule
& \multicolumn{4}{c}{แบบเชิงประจักษ์ (เท่าของ residual จริง)} & \\
\cmidrule(lr){2-5}
& 0.5$\times$ & 1$\times$ & 2$\times$ & 4$\times$ & 8$\times$ \\
Jaccard@5 & 0.968 & 0.952 & 0.935 & 0.920 & 0.912 \\
\bottomrule
\end{tabular}
\end{table}

ผลที่เด่นชัดคือ \textbf{ไม่พบจุดที่อันดับพังแม้ขยายสัญญาณรบกวนไปถึง 8 เท่าของ residual จริง}
Jaccard@5 ลดลงอย่างช้า ๆ จาก 0.979 ที่ $\pm10\%$ เหลือ 0.954 ที่ $\pm50\%$ และเหลือ 0.912
ที่ความคลาดเคลื่อนใหญ่กว่าที่วัดได้จริงถึง 8 เท่า ซึ่งสอดคล้องกับผลในหัวข้อ~\ref{sec:exp6}
ว่าอันดับ node ทนทานกว่าที่ fit error เพียงอย่างเดียวจะทำนายได้

ตารางที่~\ref{tab:noise_model} แยกตามโมเดลน้ำหนัก เพื่อตรวจว่าความทนทานเป็นคุณสมบัติของ
โมเดลใดโมเดลหนึ่งหรือของโครงสร้างกราฟเอง

\begin{table}[H]
\centering
\caption{Jaccard@5 แยกตามโมเดลน้ำหนัก (แบบสมมาตร $\pm50\%$ และเชิงประจักษ์ $8\times$)}\label{tab:noise_model}
\small
\begin{tabular}{lrrrr}
\toprule
& epss & cvss & uniform & random (null) \\
\midrule
สมมาตร $\pm50\%$ & 0.986 & 0.909 & 0.966 & 0.956 \\
เชิงประจักษ์ $8\times$ & 0.921 & 0.908 & 0.914 & 0.905 \\
\bottomrule
\end{tabular}
\end{table}

ที่สำคัญคือ \textbf{แม้แต่โมเดล \texttt{random} ซึ่งสุ่มน้ำหนักทิ้งทั้งหมด} ก็ยังให้ Jaccard@5
สูงถึง 0.905--0.956 เมื่อถูกรบกวนเพิ่ม นั่นหมายความว่าความทนทานที่สังเกตได้ไม่ได้มาจาก
คุณภาพของน้ำหนักที่ใช้ แต่มาจาก\textbf{โครงสร้างของกราฟเอง} — มีจำนวน node จำกัดที่อยู่บน
เส้นทางถูกสุดได้ และการรบกวนน้ำหนักเพียงเล็กน้อยไม่พอจะเปลี่ยนว่า node ใดอยู่ในกลุ่มนั้น
ข้อสรุปนี้ต้องอ่านคู่กับหัวขัดถัดไป ซึ่งแสดงว่าอันดับไม่ทนทานต่อ\emph{โครงสร้าง}ที่เปลี่ยน

\section{ความเสถียรของอันดับเมื่อ edge ขาดหาย}

การทดลองนี้ลบ edge แบบสุ่ม 5, 10 และ 20\% ของกราฟ เพื่อจำลองกรณีที่ข้อมูลเก็บมาไม่ครบ
ต่างจากการรบกวนน้ำหนักในหัวข้อก่อน การลบ edge เปลี่ยน\emph{โครงสร้าง}ของกราฟโดยตรง
และอาจตัดบางคู่ (entry, target) ออกจากกันทั้งคู่

\begin{table}[H]
\centering
\caption{ผลของการลบ edge แบบสุ่ม}\label{tab:drop}
\small
\begin{tabular}{rrrrp{2.6cm}}
\toprule
edge ที่ลบ & Jaccard@5 ดิบ & Kendall $\tau$ & คู่ที่ยังถึงกัน & Jaccard@5 เงื่อนไข (n กราฟ) \\
\midrule
5\%  & 0.493 & 0.672 & 69\% & 0.749 (177) \\
10\% & 0.321 & 0.526 & 47\% & 0.650 (97) \\
20\% & 0.145 & 0.371 & 20\% & 0.460 (17) \\
\bottomrule
\end{tabular}
\end{table}

คอลัมน์ ``ดิบ'' คือ Jaccard@5 เฉลี่ยทุกกราฟ ซึ่งปนผลสองอย่างเข้าด้วยกัน: อันดับเปลี่ยนเพราะ
ต้นทุนเส้นทางเปลี่ยน และคู่ (entry, target) บางคู่หลุดจากกันไปเลยทำให้ node ที่เคยอยู่ในอันดับ
หายไปโดยสิ้นเชิง เพื่อแยกสองผลนี้ คอลัมน์ ``เงื่อนไข'' คำนวณเฉพาะกราฟที่ยังคงคู่ (entry, target)
ไว้ได้อย่างน้อย 90\% หลังลบ edge แม้เงื่อนไขนี้จะทำให้ Jaccard@5 สูงขึ้นมาก (0.749 เทียบ 0.493
ที่ 5\%) แต่ก็ยังลดลงต่อเนื่องเมื่อลบมากขึ้น และจำนวนกราฟที่ผ่านเงื่อนไขลดฮวบเหลือเพียง 17 กราฟ
ที่ 20\% ทำให้ตัวเลขนั้นมีความแปรปรวนสูง (ช่วงความเชื่อมั่น 95\% กว้างถึง 0.32--0.66)

ผลสำคัญคือ \textbf{edge ที่ขาดหายกระทบอันดับรุนแรงกว่าน้ำหนักที่คลาดเคลื่อนมาก} ลบเพียง 5\%
ของ edge ลด Jaccard@5 ลงมากกว่าการรบกวนน้ำหนักถึง 8 เท่าของ residual จริงเสียอีก
(0.493 เทียบ 0.912) ข้อสรุปเชิงปฏิบัติคือ ความสมบูรณ์ของข้อมูล (การเก็บ edge ให้ครบ)
สำคัญกว่าความแม่นยำของน้ำหนักบน edge ที่เก็บมาได้แล้ว

\section{การเปรียบเทียบโมเดลน้ำหนัก 4 แบบ}

หัวข้อนี้เทียบอันดับจุดตัดที่ได้จากน้ำหนัก 4 แบบบนกราฟโครงสร้างเดียวกัน (paired comparison):
EPSS, CVSS (คะแนนฐานหารสิบ), uniform ($p=0.5$ ทุก edge ที่ไม่ใช่ deterministic) และ random
(สุ่ม $p$ ใหม่ทุกครั้ง) โมเดลหลังสุดเป็น null model: โครงสร้างใดก็ตามที่รอดจากการสุ่มทิ้งนี้
คือโครงสร้างจริง ไม่ใช่สัญญาณจากน้ำหนัก

\begin{table}[H]
\centering
\caption{Jaccard@5 ระหว่างคู่โมเดลน้ำหนัก (216 กราฟ จับคู่กัน)}\label{tab:models}
\small
\begin{tabular}{lrr}
\toprule
คู่โมเดล & Jaccard@5 & Kendall $\tau$ \\
\midrule
epss -- uniform & 0.974 & 0.947 \\
epss -- cvss & 0.951 & 0.884 \\
epss -- random (null) & 0.924 & 0.834 \\
uniform -- random (null) & 0.929 & 0.832 \\
cvss -- uniform & 0.941 & 0.878 \\
cvss -- random (null) & 0.910 & 0.817 \\
\bottomrule
\end{tabular}
\end{table}

ตัวเลขที่ต้องอ่านคู่กันคือแถวที่สองและสาม: EPSS เทียบกับ CVSS ให้ Jaccard@5 เท่ากับ 0.951
ขณะที่ EPSS เทียบกับ \emph{random} ให้ 0.924 — ต่างกันเพียง 0.027 \textbf{จึงไม่มีหลักฐานเพียงพอ
ที่จะอ้างว่า EPSS ให้อันดับที่ ``ถูกต้องกว่า'' CVSS} เพราะกราฟจำลองไม่มีค่าจริง (ground truth)
ให้เทียบ ตัวเลขทั้งหมดในตารางนี้วัดได้แค่ \emph{ความต่างกันของอันดับ} ระหว่างสองน้ำหนัก
ไม่ใช่ความถูกต้อง การที่ EPSS ใกล้ CVSS มากกว่าใกล้ random เพียงเล็กน้อย บ่งชี้ว่าทั้งสองโมเดล
ที่มีข้อมูลจริงรองรับให้อันดับที่คล้ายกันในระดับหนึ่ง แต่ความคล้ายนั้นเทียบกับ null model แล้ว
มีขนาดเล็ก สรุปว่าการทดลองนี้วัดได้เพียง disagreement rate และ
ความเสถียรของอันดับ ไม่ใช่ความแม่นยำ
\section{ผลของความคลาดเคลื่อนจากการ fit ต่ออันดับจุดตัด}\label{sec:exp6}\draftnote{§5.7 — exp 6}
\section{สัดส่วนเส้นทางที่ข้ามโดเมน}\draftnote{§5.8 — exp 7}
\section{ผลจาก lab จริง (วัตถุประสงค์ข้อ 5)}\draftnote{§5.9}
\section{อภิปรายผลและภัยต่อความตรงของผล}\draftnote{§5.10}
\section{ข้อจำกัด}\draftnote{§5.11}
